% Figure: ridge and lasso coefficient paths for four coefficients as
% the penalty strength grows. Ridge shrinks every coefficient smoothly
% toward zero without reaching it; lasso drives coefficients to exactly
% zero one after another, which is its variable-selection property. The
% x-axis is increasing penalty (log scale suggested but drawn linearly
% in a coded penalty). Coordinates are schematic paths that reproduce
% the qualitative contrast the section describes.
\begin{figure}[htbp]
\centering
\begin{tikzpicture}
\begin{axis}[
    name=ridgeax,
    width=7.6cm, height=6.6cm,
    xmin=0, xmax=10,
    ymin=-1.2, ymax=2.6,
    axis lines=left,
    title={\small Ridge ($\ell_2$ penalty)},
    xlabel={Penalty $\lambda$},
    ylabel={Coefficient $\hat{\beta}_j$},
    every axis plot/.append style={very thick, mark=none},
    xtick={0,2,4,6,8,10},
]
% Ridge: each coefficient decays as b0/(1+c*lambda), never reaching 0.
\addplot[engblue,   domain=0:10, samples=40]{ 2.4/(1+0.55*x) };
\addplot[engred,    domain=0:10, samples=40]{ 1.5/(1+0.55*x) };
\addplot[englime!70!black, domain=0:10, samples=40]{ 0.8/(1+0.55*x) };
\addplot[engamber,  domain=0:10, samples=40]{ -1.0/(1+0.55*x) };
\addplot[enggray!50, dashed] coordinates {(0,0) (10,0)};
\end{axis}
\begin{axis}[
    at={(ridgeax.south east)}, anchor=south west, xshift=1.2cm,
    width=7.6cm, height=6.6cm,
    xmin=0, xmax=10,
    ymin=-1.2, ymax=2.6,
    axis lines=left,
    title={\small Lasso ($\ell_1$ penalty)},
    xlabel={Penalty $\lambda$},
    ylabel={Coefficient $\hat{\beta}_j$},
    every axis plot/.append style={very thick, mark=none},
    xtick={0,2,4,6,8,10},
]
% Lasso: piecewise-linear decay, each hitting exactly zero and staying.
% beta1 (blue): from 2.4, hits 0 at lambda=9.
\addplot[engblue] coordinates {(0,2.4) (9,0) (10,0)};
% beta2 (red): from 1.5, hits 0 at lambda=6.
\addplot[engred] coordinates {(0,1.5) (6,0) (10,0)};
% beta3 (lime): from 0.8, hits 0 at lambda=3.
\addplot[englime!70!black] coordinates {(0,0.8) (3,0) (10,0)};
% beta4 (amber): from -1.0, hits 0 at lambda=4.
\addplot[engamber] coordinates {(0,-1.0) (4,0) (10,0)};
\addplot[enggray!50, dashed] coordinates {(0,0) (10,0)};
\end{axis}
\end{tikzpicture}
\caption[Ridge and lasso coefficient paths]{Coefficient paths for four
predictors as the penalty strength grows. Ridge, with its squared
$\ell_2$ penalty, shrinks every coefficient smoothly toward zero and
never reaches it, so all four predictors stay in the model at every
penalty. Lasso, with its $\ell_1$ penalty, drives coefficients to
exactly zero one after another, so the model selects variables as the
penalty rises: at a moderate penalty only the largest coefficients
survive. The contrast is the geometric consequence of the penalty
shape, the $\ell_1$ ball's corners sitting on the coordinate axes,
and it is the reason lasso is the reach when many candidate predictors
are believed inactive and ridge is the reach when all predictors are
kept but the design is ill-conditioned.}
\label{fig:vol02:ch14:shrinkage-paths}
\end{figure}
